\Subhead{2.4.2}{Сплошная среда со слабыми периодическими неоднородностями. Связанные волны}

В периодической среде решение имеет вид $u_k(x)=\e^{\ii kx}p_k(x)$, $p_k(x+a)=p_k(x)$. При слабой модуляции сильнее всего связываются волны, чьи волновые числа отличаются на вектор обратной решётки $G=2\pi/a$. Вблизи $2k\simeq G$ достаточно встречных амплитуд $A,B$. Выбирая фазу так, чтобы коэффициент связи $\kappa$ был вещественным, получаем пространственную модель
\begin{equation}
 A'=\ii\delta A+\ii\kappa B,\qquad
 B'=-\ii\delta B-\ii\kappa A.
\end{equation}
Для $A,B\propto\e^{\ii qx}$ имеем $q^2=\delta^2-\kappa^2$. При $|\delta|<|\kappa|$ $q$ мнимо: в бесконечной решётке нет распространяющейся блоховской моды, поле проникает лишь на конечную глубину. Это запрещённая полоса, возникающая из многократного когерентного отражения. В конечной структуре наблюдается сильное отражение, а не обязательное поглощение энергии. Устойчивость во времени и пространственное затухание в запрещённой полосе — разные понятия.
